\section{Quantities in Types}

\label{sec:quantities}

The biggest distinction between Idris 2 and its predecessor is that it is
built on Quantitative Type Theory (QTT)~\cite{nuttin,qtt}. In QTT, each
variable binding (including function arguments) has a \emph{multiplicity}.
QTT itself takes a general approach to multiplicities, allowing any semiring.
For Idris 2, we make a concrete choice of semiring, where a
multiplicity can be one of:

\begin{itemize}
\item $0$: the variable is not used at run time
\item $1$: the variable is used exactly once at run time
\item $\omega$: no restrictions on the variable's usage at run time
\end{itemize}

In this section, we describe the syntax and informal semantics of
multiplicities, and discuss the most important practical applications:
erasure, and linearity. The formal semantics are presented
elsewhere~\cite{qtt}; here, we aim to describe the intuition.
In summary:

\begin{itemize}
\item Variable bindings have multiplicities which describe how often the
varible must be used within the scope of its binding.
\item A variable is ``used'' when it appears in the body of a definition (that
is, not a type declaration), in an argument position with multiplicity $1$ or
$\omega$.
\item A function's type gives the multiplicities the arguments have in the
function's body.
\end{itemize}

Variables with multiplicity $\omega$ are truly unrestricted, meaning that they
can be passed to argument positions with multiplicity $0$, $1$ or $\omega$.
A function which
takes an argument with multiplicity $1$ promises that it will not share the
argument in the future; there is no requirement that it has not been shared
in the past.

\subsection{Syntax}

When declaring a funcion type we can, optionally, give an explicit multiplicity
of \texttt{0} or \texttt{1}. If an argument has no multiplicity
given, it defaults to $\omega$.  For example, we can declare the type of
an identity function which takes its polymorpic type \emph{explicitly}, and
mark it erased:

\begin{lstlisting}
id_explicit : (0 a : Type) -> a -> a 
\end{lstlisting}

If we give a partial definition of \FN{id\_explicit}, then inspect the type of
the hole in the right hand side, we can see the multiplicities of the variables
in scope:

\begin{lstlisting}
id_explicit a x = ?id_explicit_rhs

Main> :t id_explicit_rhs
0 a : Type
  x : a
------------------------------
id_explicit_rhs : a
\end{lstlisting}

This means that \VV{a} is not available at run time, and \VV{x} is unrestricted
at run time. If there is no explicit multiplicity given, it is $\omega$. A
variable which is not available at run time can only be used in
argument positions with multiplicity \texttt{0}.
%
Implicitly bound arguments are also given
multiplicity \texttt{0}. So, in the following declaration of \FN{append}\ldots

\begin{lstlisting}
append : Vect n a -> Vect m a -> Vect (n + m) a
\end{lstlisting}

\ldots \FN{n}, \FN{a} and \FN{m} have multiplicity \texttt{0}. In Idris 2,
therefore, unlike Idris 1, the declaration is equivalent to writing:

\begin{lstlisting}
append : {0 n : Nat} -> {0 m : Nat} -> {0 a : Type} ->
         Vect n a -> Vect m a -> Vect (n + m) a
\end{lstlisting}

The default multiplicities for arguments are, therefore:

\begin{itemize}
\item If you explicitly write the binding, the default is $\omega$.
\item If you omit the binding (e.g. in a type level variable), the default is $0$.
\end{itemize}

This design choice, that type level variables are by default
erased, is primarily influenced by the common usage in Idris 1, that
implicit type level variables are typically compile-time only.
As a result, the most common use cases involve the most lightweight syntax.

\subsection{Erasure}

\label{sec:erasure}

The multiplicity \texttt{0} gives us control over whether a function
argument---\TC{Type} or otherwise---is used at run time. This is important in
a language with first-class types since we often parameterise types by values
in order to make relationships between data explicit, or to make assumptions
explicit. In this section, we will consider two examples of this, and see
how multiplicity 0 allows us to control which data is available at run time.

\subsubsection{Example 1: Vector length}

We have seen how a vector's type includes its length;
we can use this length at run time, even though it is
part of the type, provided that it has non-zero multiplicity:

\begin{lstlisting}
length : {n : Nat} -> Vect n a -> Nat
length xs = n
\end{lstlisting}

With this definition, the length \VV{n} is available at run time, which has
a storage cost, as well as potentially the cost of computing the length
elsewhere. If we want the length to be erased, we would need to recompute it
in the \FN{length} function:

\begin{lstlisting}
length : Vect n a -> Nat
length [] = Z
length (_ :: xs) = S (length xs)
\end{lstlisting}

The type of \FN{length} in each case makes it explicit whether or not the
length of the \TC{Vect} is available. Let us now consider a more realistic
example, using the type system to ensure soundness of a compressed encoding
of a list.

\subsubsection{Example 2: Run-length Encoding of Lists}

Run-length encoding is a compression scheme which collapses sequences (runs)
of the same element in a list. It was originally developed for reducing the
bandwidth required for television pictures~\cite{rle}, and later used in
image compression and fax formats, among other things.

We will define a data type for storing run-length encoded lists, and use the
type system to ensure that the encoding is sound with respect to the
original list. To begin, we define a function which constructs a list by
repeating an element a given number of times. We will need this for
explaining the relationship between compressed and uncompressed data.

\begin{lstlisting}
rep : Nat -> a -> List a
rep Z x = []
rep (S k) x = x :: rep k x
\end{lstlisting}

Using this, and the concatenation operator for \TC{List} (\FN{++}, which is
defined like \FN{append}), we can describe what it means to be a run-length
encoded list:

\begin{lstlisting}
data RunLength : List ty -> Type where
     Empty : RunLength []
     Run : (n : Nat) -> (x : ty) -> (rle : RunLength more) ->
           RunLength (rep (S n) x ++ more)
\end{lstlisting}

%We can read this as:

\begin{itemize}
\item \DC{Empty} is the run-length encoding of the empty list \DC{[]}
\item Given a length \VV{n}, an element \VV{x}, and an encoding \VV{rle} of the
list \VV{more}, \DC{Run} \VV{n} \VV{x} \VV{rle} is the encoding
of the list \texttt{rep (S n) x ++ more}. That is, $n+1$ copies of \texttt{x}
followed by \VV{more}.
\end{itemize}

We use \texttt{S n} to ensure that \texttt{Run} always increases the length
of the list, but otherwise we make no attempt (in the type) to ensure that
the encoding is optimal; we merely ensure that the encoding is sound with
respect to the encoded list.
%
Let us try to write a function which uncompresses a run-length encoded list,
beginning with a partial definition:

\begin{lstlisting}
uncompress : RunLength {ty} xs -> List ty
uncompress rle = ?uncompress_rhs
\end{lstlisting}

\textbf{Note:} The \VV{\{ty\}} syntax gives an explicit value for the implicit
argument to \TC{RunLength}. This is to state in the type that the element type
of the list being returned is the same as the element type in the encoded
list.

Like our initial implementation of \FN{length} on \TC{Vect}, we might be
tempted to return \VV{xs} directly, since the index of the encoded list gives
us the original uncompressed list. However, if we check the type of
\texttt{uncompress\_rhs}, we see that \texttt{xs} isn't available at run time:

\begin{lstlisting}
 0 xs : List ty
   rle : RunLength xs
------------------------------
uncompress_rhs : List ty
\end{lstlisting}

This is a good thing: if the uncompressed list were available at
run-time, there would have been no point in compressing it!
We can still take advantage of having the uncompressed list available
as part of the type, though, by refining the type of \FN{uncompress}:

\begin{lstlisting}
data Singleton : a -> Type where
     Val : (x : a) -> Singleton x

uncompress : RunLength xs -> Singleton xs
\end{lstlisting}

A value of type \texttt{Singleton x} has exactly one value, \texttt{Val x}.
The type of \FN{uncompress} expresses that the uncompressed list
is guaranteed to have the same value as the original list, although it must
still be reconstructed at run-time. We can implement it as follows:

\begin{lstlisting}
uncompress : RunLength xs -> Singleton xs
uncompress Empty = Val []
uncompress (Run n x y) = let Val ys = uncompress y in
                             Val (x :: (rep n x ++ ys))
\end{lstlisting}

\textbf{Aside:} This implementation was generated by type-directed program
synthesis~\cite{polikarpova_program_2016}, rather than written by hand,
taking advantage of the explicit relationship given in the type between the
input and the output. The definition of \TC{RunLength} more or less directly
gives the uncompression algorithm, so we should not need to write it again!

The \texttt{0} multiplicity, therefore, allows us to reason about values
at compile time, without requiring them to be present at run time.
Furthermore, QTT gives us a guarantee of erasure, as well as an
explicit type level choice as to whether an index is erased or not.
